\begin{solapartat}
\punts{0,2} \[ \left|\vec{E}\right| = \frac{\Delta V}{d} = \frac{220}{0,2} = 1100\ \mathrm{V/m} \]

\punts{0,2} \[ \vec{F}_E = q\vec{E} \]

\figura[0.08]{sol_fig1.png}

\punts{0,3} Per que l'electró es mogui en línia recta vertical i cap amunt, la força aplicada pel camp elèctric ha de ser vertical i sentit cap amunt (dirigida del càtode cap a l'ànode).

La força i el camp elèctric són paral·lels, llavors la direcció del camp elèctric és vertical, perpendicular a les superfícies de l'ànode i del càtode.

\punts{0,3} Com la càrrega de l'electró és negativa el sentit del camp elèctric és l'oposat al de la força, és a dir, va dirigit cap avall, de l'ànode cap al càtode.

Així les línies de camp elèctric són:

\figura[0.25]{sol_fig2.png}

\punts{0,25} Les línies de camp indiquen la direcció en la que el potencial disminueix, per tant, el potencial ha de ser més gran a l'ànode que al càtode.
\end{solapartat}

\begin{solapartat}
Velocitat quan surten de l'ànode:

\destaca{Primera opció, per forces}

\punts{0,2} \[ a = \frac{F}{m} = \frac{eE}{m} = 1,93\times 10^{14}\ \mathrm{m/s^2}. \]

\punts{0,2} \[ v_f^2 - v_i^2 = 2a\Delta s \Rightarrow v = \sqrt{2a\Delta s} = \sqrt{2\times 1,93\times 10^{14}\times 0,2} = 8,80\times 10^{6}\ \mathrm{m/s} \]

\destaca{Alternativament, per energies:}

\punts{0,3} \[ \tfrac{1}{2} m v_f^2 - \tfrac{1}{2} m v_i^2 = W_{Tot} = -q\Delta V = e\Delta V = 3,52\times 10^{-17}\ \mathrm{J}. \]

\punts{0,1} \[ v = \sqrt{2\,\frac{e\Delta V}{m}} = 8,80\times 10^{6}\ \mathrm{m/s} \]

\punts{0,3} Velocitat quan impacten en el blanc: el camp magnètic aplica una força normal a la trajectòria, de manera que no produeix cap canvi en el mòdul de la velocitat, només provoca un canvi en la direcció del moviment, per tant, $v = 8,80\times 10^{6}\ \mathrm{m/s}$.

\punts{0,2} Per tal que els electrons descriguin la trajectòria circular del dibuix, l'acceleració normal o centrípeta ha de tenir el sentit indicat a la figura. Llavors, com la força i l'acceleració tenen la mateixa direcció i sentit, la força magnètica ha d'anar dirigida segons l'eix $y$, com s'indica a la figura.

\figura[0.4]{sol_fig3.png}
\figura[0.4]{sol_fig4.png}

\punts{0,1} Sabem que $\vec{F}_m = q\vec{v}\times\vec{B}$, llavors $\vec{B}$ ha de ser perpendicular a $\vec{v}$ i $\vec{F}_m$, per tant, segons els eixos del dibuix, $\vec{B}$ ha d'anar dirigit segons l'eix $x$.

\punts{0,25} El sentit el podem establir a partir de la regla de la mà dreta, perquè el producte vectorial $\vec{v}\times\vec{B}$ estigui dirigit segons l'eix $y$ en el sentit de la figura, $\vec{B}$ hauria de tenir el sentit positiu segons l'eix de les $x$, però com la càrrega de l'electró és negativa, el sentit final és l'oposat, per tant $\vec{B}$ ha de tenir el sentit negatiu segons l'eix de les $x$. Per tant, representem $\vec{B}$ perpendicular al pla del dibuix i sentit cap endins, és a dir, el representem amb creus.
\end{solapartat}
